\documentclass[../main.tex]{subfiles}
\begin{document}
\chapter{Costos Variables}
\img{images/05/apple-fabrica.jpg}{Fábrica de Apple en China.}
\info{Costo Variable}{El costo variable es el gasto que fluctúa en
  proporción a la actividad generada por una empresa o, en otros
  términos, el que depende de las variaciones que afecten a su
  volumen de negocio.
}
Se puede dar el caso de que si una organización se dedica a la
producción de vino -una bodega- necesitará como materia prima una
buena cosecha de uva de tal modo que si incrementa sus índices de
actividad, requerirá mayor cantidad de producto y, como consecuencia,
también verá aumentados sus costos variables. Como se puede apreciar
en el gráfico, con incrementos de producción -volumen- se producen
incrementos de costos variables.
\section{Tipos de costo variable}
Como su propia naturaleza indica, los costos variables cambian
atendiendo al número de unidades producidas en una organización, en
relación a su volumen de negocio. Por este motivo, puede clasificarse
en tres categorías diferentes:
\begin{itemize}
  \item costo variable proporcional: Se corresponde con el que varía
    en la misma proporción que el nivel de producción de la entidad;
    asimismo, el costo variable unitario se mantiene constante.

  \item costo variable progresivo: Se relaciona con el que cambia más
    que proporcionalmente ante variaciones del nivel de producción;
    por su lado, el costo variable unitario es creciente.

  \item  costo variable regresivo: Define el que fluctúa menos que
    proporcionalmente a variaciones en el nivel de producción. El
    costo variable unitario es decreciente.
\end{itemize}
\section{Características del costo variable}
\begin{center}
  \begin{tikzpicture}
    \draw (0,0)--(0,5)[-stealth] node [left] {C:costo};
    \draw (0,0)--(5,0)[-stealth] node [below] {Q:cantidad producida};
    \draw (0,0)--(5,5) node [below right] {CV:Costos Variables};
  \end{tikzpicture}
\end{center}
Las características del costo variable son:
\begin{enumerate}
  \item Si la producción de artículos, bienes o servicios se anula,
    los costos variables desaparecen.
  \item La cantidad de costos variables tenderán a ser proporcionales
    a la cantidad de bienes producidos.
  \item Los costos variables no dependen del tiempo sino, como ya se
    ha subrayado, del volumen de negocio de la empresa.
  \item Este tipo de gasto se puede controlar y gestionar a corto plazo.
  \item Está regulado y clasificado por el departamento de
    administración de la entidad.
\end{enumerate}
En resumidas cuentas, el costo variable puede ayudar a comprobar los
resultados económicos de una organización ofreciendo una información
exacta del comportamiento del negocio: si la actividad de producción
aumenta, este tipo de gasto también se incrementará y, viceversa, si
aquélla disminuye o cae, el costo variable responderá de modo similar.
\actividad{ Actividad para el proyecto}
\begin{enumerate}
  \item Identificar los costos variables.
  \item Identificar como impacta la cantidad de producción en los
    costos variables, determinar el costo unitario variable, y el
    costo unitario total.
\end{enumerate}
\actividad{ Investigar y Contestar}
\begin{enumerate}
  \item ¿Cuál es la diferencia entre costo y gasto?
  \item ¿Costos variables que pueden acumularse, se transforman en
    fijos? Justifica la respuesta.
  \item Definir el costo unitario.
  \item Definir el costo unitario variable.
  \item Definir el costo total.
\end{enumerate}

\end{document}
